\section{A graph cut polynomial in the second block}
The preceding algorithm becomes expensive when many variables contain
no whole clause among them. We give a family for which the first
block has a known value and evaluating the second is already
\(\#\mathrm P\)-hard. Precision two here means two blocks of \(s\)
base-13 digits. The reduction uses polynomially growing \(s\).
The cut polynomial in the theorem is the zero-field Ising partition
function up to a scalar. Indeed, with
\(\epsilon_v=(-1)^{\sigma(v)}\) and \(k=|E(G)|\),
\(2|\delta(\sigma)|=k-\sum_{uv\in E(G)}\epsilon_u\epsilon_v\).
For positive \(Y>1\), the edge weights favor unequal spins.
For the conventional Ising formulation see \cite{js}. The construction
below places this graph polynomial at a particular arithmetic position
of a sparse univariate root count.

\begin{theorem}\label{th:cuts}
Let \(G\) have \(q\) vertices, no loops or isolated vertices, and \(k\)
edges. Write \(Z_G(Y)=\sum_{\sigma:V(G)\to\{0,1\}}Y^{|\delta(\sigma)|}\),
where \(\delta(\sigma)\) is the set of edges whose endpoints differ.
Given \(G\) and a positive integer \(h\) in unary, one can construct
a normalized formula in which every row contains the same positive
literal \(x\), and \(x\) is forced true. The formula has exactly
\(2^{q+2}\) witnesses. From \(R_2\) of its field polynomial one can
recover \(Z_G(2^h)\bmod 13^s\) in polynomial bit time.
With \(h=q+1\) and a polynomial choice of \(s\), a single evaluation
of its second block recovers every coefficient of \(Z_G\).
\end{theorem}

\begin{proof}
Use \(N=q+3\) variables: the vertices of \(G\), the center \(x\),
and two additional bits \(y,z\). For each edge \(uv\) put in \(h\)
copies of each ordered row \((x,u,v),(x,v,u)\).
Append four binary rows \((x,t,t)\) for \(t=y,\neg y,z,\neg z\).
There are \(g=2hk\) edge rows and \(M=g+4\) rows in total. Two
opposite requirements on \(y\) force \(x\) true; then every other
bit is free. The second pair of opposite requirements will also
exclude false-center contributions at first order.

Introduce a formal variable \(t\) for an equal edge with unequal-edge
weight one. After removing the four required equal edges from each
gate, let \(f_{pqr}(t)=c_0+c_1t+\cdots\) be its local polynomial.
More precisely, replace each \(w\) in \eqref{eq:local} by \(t\)
on equal bits and 1 otherwise, and divide by \(t^4\). The local gate
has 25 edges, so \(\widehat K=\beta^{21} f_{pqr}(\alpha/\beta)\).
Writing \(b(\sigma)\) for the number of defective variable pairs,
form the integer polynomial
\[
 P(t)=\sum_{\sigma\in\{00,01,10,11\}^{N}}
            t^{b(\sigma)}\prod_{c}f_c(\sigma;t).
\]
In \(\mathbb Z/\alpha^2\mathbb Z\), let \(\tau=\alpha\beta^{-1}\)
and \(T=N+21M+1\). The vertex and anchor factors give
\begin{equation}
 V=\frac{R_2}{\beta^T(1+\tau)}
   =P(0)+\tau P'(0)\pmod{\alpha^2},\qquad P(0)=2^{N-1}.       \label{eq:derivative}
\end{equation}
All denominators in this equality are units. Subtract \(2^{N-1}\)
from a representative of \(V\), divide by \(\alpha\), and multiply
by \(\beta\). The result is \(P'(0)\bmod\alpha\).

Here are the entries of the local gate needed for the derivative.
Use \(U=(0,1)\) for true and \(F=(1,0)\) for false. At \(p=U\)
the leading coefficient is one for \emph{all} states of \(q',r\),
including defects. On healthy \(q',r\) the first correction is
\[
 \begin{array}{c|rr}
 c_1(U,q',r)&r=U&r=F\\\hline
 q'=U&8&8\\ q'=F&12&9
 \end{array}.
\]
For a verification by hand, the triangle polynomial
\(L(d_1,d_2,d_3)/t\bmod t^2\) with these weights is \(3t\)
when the three neighboring bits coincide and \(1+2t\) otherwise.
Indeed a mixed triple forces a unique internal choice with one equal
edge, and two choices give two equal edges. For three equal bits
three choices give two equal edges. Multiplying these four linear
polynomials and the private-edge weight in \eqref{eq:local} gives
the displayed table. At constant order the private pair \(u,v\)
must be complementary. Its choice \(v=1\) works just when
\(p_-=1\) or \(q'_-=1\). The choice \(v=0\) works just when
\(p_+=q'_+=r_-=1\) and at least one of \(p_-,q'_-\) is zero.
Consequently on healthy \(q',r\),
\begin{align*}
 c_0(11,q',r)&=1+\mathbf1_{\{q'=F,\ r=U\}},\\
 c_0(00,q',r)&=\mathbf1_{\{q'=U\}} .
\end{align*}

Only zero or one defective pair can contribute to \(P'(0)\).
With no defect, a false center violates at least one pin row for
each of \(y,z\), giving a double zero. For a true center all leading
entries equal one. The product derivative adds the local corrections.
Their sum over all assignments on two distinct neighboring variables
is \(37\), and on one repeated signed variable it is \(17\).
Thus all such derivatives together give
\(37g\,2^{N-3}+17\cdot4\,2^{N-2}\).

If any variable other than \(x\) is the sole defect, a false center
still has a healthy contradictory pair of pin rows. For a true
center all local leading coefficients are one. These possibilities
contribute \((N-1)2^{N-1}\). For a sole defect \(x=00\), the two
signs on \(y\) cannot both be true and give zero. For \(x=11\)
each pin row gives one. An edge \(uv\) gets a factor two in exactly
one of the two input orders if its ends differ, and gets one in both
orders if they agree. The two independent pin bits contribute four.
Putting the contributions together gives the integer identity
\begin{equation}
 P'(0)=37g\,2^{N-3}+68\,2^{N-2}+(N-1)2^{N-1}
                           +4Z_G(2^h).                              \label{eq:cut}
\end{equation}
Four is a unit modulo \(\alpha\). Equations \eqref{eq:derivative}
and \eqref{eq:cut} recover the claimed value.

Each coefficient of \(Z_G\) is at most \(2^q\). Choose \(h=q+1\)
and \(13^s>2^{q+hk}\), with \(s\ge2\). Then \(Z_G(2^h)\) is the
actual integer and its base-\(2^h\) digits are the coefficients
without carries. All field parameters, the bit lengths of weights
and integers, and the number of written terms in \(H\) are polynomial
in \(q+k\) by the sparse construction \cite[Thm.~1]{degree}.
The lower block is already known from \eqref{eq:derivative}, so only
the second block needs an oracle.
\end{proof}
